\section{Structural relations and further questions}\label{intro:discussion}

Both models begin with the full collision difference, but resonance
compatibility enters at different stages. In the pinned chain, it
relates values at different points, gives regularity of measurable
solutions, and determines all invariants through periodic matching.
At critical resonances, the same geometry has another consequence:
the energy mismatch and the function difference share a vanishing
factor, allowing the singular measure and the squared difference to
be treated together. Kernel classification, compact resonance
averaging and construction of the closed form thus form a connected
argument.

In the three-dimensional model, collision invariants determine both
the equilibrium family and the error decomposition. Actual five-moment
matching cancels the cross term in relative entropy. Collisions control
deviations that leave these moments unchanged, while transport makes the
moments vary in space. The residual flux joins these effects in the
microscopic entropy balance \eqref{cur:entropyidentity}. That balance
determines the strong resonance-current limit and then the first-order
moment fluxes. The Onsager tensor defined by the weighted cell problem
therefore also appears in the limit of microscopic entropy production
and in the local macroscopic entropy law.

Corner memory specifies the topology of this macroscopic description.
The free-transport trace \eqref{intro:corner-memory} persists throughout
the existence interval, while the distribution approaches the Euler
equilibrium in momentum-integrated norms. The estimates control the
contribution of the corner layer to total dissipation and flux;
pointwise relaxation at the corners does not hold. The first-order
microscopic quantities and resonance currents converge in weighted,
time-integrated spaces. Their strong traces at the initial time
require an initial-layer analysis.

Several further questions arise. For the pinned chain, a spatial
diffusion limit requires simultaneous rescaling of transport and
collisions, together with flux and nonlinear stability estimates
for the two conserved densities. The nondegeneracy constants of the
resonance charts change near the acoustic endpoint $d=1/2$ and the
zero-coupling endpoint $d=0$; uniform estimates there need separate
analysis.

For the three-dimensional model, the first-order law describes the
five-moment parameters of the actual kinetic solution. Comparing them
with an independently solved nonlinear viscous macroscopic system
requires stability of that system and initial-layer estimates.
Enlarging the momentum domain requires control of the equilibrium,
collision frequency and transport coefficients as the cutoff radius
$R$ grows. Finally, connecting this strong-collision limit to
inhomogeneous Schr\"odinger dynamics requires a kinetic derivation
matched to local RJ initial data, with uniform control of the weak
nonlinearity, large-volume and strong-collision parameters.

\paragraph{Companion code.}
The manuscript and Lean sources are available in the companion repository:
\begin{center}
\url{https://github.com/veridiscoverLab/resonance-kinetic-lean}.
\end{center}
The repository documents the theorem statements, their formal entry
points and the verification procedure. The pinned-chain classification,
spectral gap and closed realization, the Onsager kernel and rank, and
the constant-parameter five-moment decay have completed end-to-end
formalizations. Formalization of the nonlinear limit over the prescribed
Euler interval and the final strong-current theorem remains in progress;
their paper proofs are given above.
